\documentclass[11pt]{article}
\usepackage[a4paper,margin=24mm]{geometry}
\usepackage{amsmath,amssymb,amsthm}
\usepackage[hidelinks]{hyperref}
\newtheorem{theorem}{Theorem}
\begin{document}
\title{The printed divisibility assertion in Conjecture 00000000261 fails at $d=5$}
\author{Chengzhe Gao}\date{}\maketitle

\paragraph{Statement being addressed.}
Both source languages take a fundamental discriminant $d\equiv1\pmod4$
and a fundamental Pell solution $(t+u\sqrt d)/2$, and assert
$u\nmid t$. The displayed divisor and dividend are essential.
The conventional AAC formulation concerns divisibility of the unit
coefficient by the prime discriminant, $p\nmid u$; see~\cite{FM}.
Our counterexample concerns the printed condition $u\nmid t$ only.

For $d\equiv1\pmod4$, the standard half-integral Pell presentation
uses positive integers $t,u$ of the same parity with
$t^2-du^2=\pm4$. Its least value $(t+u\sqrt d)/2>1$ is the fundamental
unit. If ``Pell solution'' instead restricts the norm to $+1$,
use $t^2-du^2=4$ and take the least value in that class.
Both conventions give a counterexample here.

\begin{theorem}
The integer $5$ is a fundamental discriminant congruent to $1$ modulo $4$.
The unique fundamental positive pair for $t^2-5u^2=\pm4$ is $(1,1)$.
For $t^2-5u^2=4$ it is $(3,1)$. In particular $u\mid t$ in both cases.
\end{theorem}
\begin{proof}
The integer $5$ is squarefree and $5\equiv1\pmod4$, so it is a
fundamental discriminant. The pair $(1,1)$ has equal parity and
$1^2-5\cdot1^2=-4$. Every positive integer pair $(t,u)$ has $t\geq1$
and $u\geq1$. Therefore
\[
 \frac{t+u\sqrt5}{2}\geq\frac{1+\sqrt5}{2},
\]
with equality only at $(t,u)=(1,1)$. This proves genuine minimality
and uniqueness among all positive Pell pairs, with no search bound.
The coefficient $u=1$ divides $t=1$.

Under the norm-$+1$ convention, $(3,1)$ has equal parity and
$3^2-5\cdot1^2=4$. Every positive solution satisfies
$t^2=5u^2+4\geq9$, hence $t\geq3$ and $u\geq1$. Consequently
\[
 \frac{t+u\sqrt5}{2}\geq\frac{3+\sqrt5}{2},
\]
with equality only at $(3,1)$. Here $u=1$ divides $t=3$.
This disproves the universal condition as printed under both conventions.
\end{proof}

\paragraph{Why positive coefficients exhaust the unit candidates.}
For the usual ring of integers $\mathbb Z[(1+\sqrt5)/2]$, write a unit
$\alpha>1$ and its conjugate as
$\alpha=(t+u\sqrt5)/2$ and $\alpha'=(t-u\sqrt5)/2$.
Its norm is $\pm1$, so $\alpha'=\pm1/\alpha$ and $|\alpha'|<1$.
Thus $t=\alpha+\alpha'>0$ and $u\sqrt5=\alpha-\alpha'>0$.
Its integral coordinates have the same parity and satisfy the
displayed Pell equation. Conversely those pairs have norm $\pm1$
and inverses in the same ring. Hence the positive-Pell minimum used
above is the usual fundamental unit, not a smaller search class.

\paragraph{Exact order representation in Lean.}
To avoid approximating $\sqrt d$, the formal project compares
nonnegative coefficient pairs by exact integer arithmetic.
For $u\leq v$, the inequality
$t+u\sqrt d\leq s+v\sqrt d$ is equivalent to
\[
 t\leq s\quad\text{or}\quad (t-s)^2\leq d(v-u)^2,
\]
where differences in the implementation are natural-number subtraction.
If $t>s$, this is obtained by squaring the nonnegative sides of
$t-s\leq(v-u)\sqrt d$; if $t\leq s$, the comparison already holds.
For $u>v$, the equivalent condition is
\[
 t\leq s\quad\text{and}\quad d(u-v)^2\leq(s-t)^2.
\]
These equivalences justify the project predicate \texttt{ValueLE}
as the exact order of the algebraic values, including their common
denominator $2$. The project proves coordinatewise monotonicity
and the equality case at a coordinatewise minimum for arbitrary
coefficients. The real-surds interpretation above is the elementary
semantic bridge; no real-number library or pre-assumed Pell
minimality is used.

\paragraph{Formal verification and limits.}
Lean defines squarefreeness over every potential squared divisor,
the appropriate fundamental-discriminant class, both Pell equations
over all positive integer pairs, and fundamentality by minimality
against every candidate. It proves the two minima, uniqueness,
divisibility and negations of the source's universal claim.
The finite checks establish only the concrete pair equations and
the small discriminant's squarefreeness; minimality is an unbounded proof.

The project uses Lean 4.19.0 with bundled Std only. It has passed a fresh
\texttt{lake build} with warnings as errors and
\texttt{lake env lean -DwarningAsError=true Main.lean}.
The printed axiom dependencies are standard foundational axioms;
there is no \texttt{sorry}, \texttt{admit}, \texttt{native\_decide}
or custom axiom. The Python script is only a supplementary exact
arithmetic check. The manuscript and proof were reviewed by the
same assistant in a separate adversarial pass, without a subagent
or claim of independent review.

This submission makes no new claim about the conventional AAC assertion
$p\nmid u$. At $d=5$ the coefficient $u=1$ does satisfy that different
condition. The disproof applies to the displayed source's $u\nmid t$.

\begin{thebibliography}{1}
\bibitem{FM} N. Fellini and M. Ram Murty,
\emph{Wieferich primes in number fields and the conjectures of
Ankeny--Artin--Chowla and Mordell}, 2025.
\href{https://arxiv.org/abs/2508.08472}{arXiv:2508.08472}.
\end{thebibliography}
\end{document}
